\section{R ile parametrik olmayan testler}
\label{sec:r-nonparametric}

\textbf{Soru.} İki bağımsız grup için Mann--Whitney, aynı kişilerin iki
ölçümü için Wilcoxon işaretli sıra ve üç bağımsız grup için Kruskal--Wallis
testlerini R ile uygulayın.

\begin{lstlisting}[style=prismR]
# Iki bagimsiz grup
a <- c(7, 8, 9, 10, 11)
b <- c(1, 2, 3, 4, 5)
wilcox.test(a, b, alternative = "two.sided",
            exact = TRUE, correct = FALSE)

# Ayni kisilerin iki olcumu
once <- c(10, 12, 15, 18, 20, 22)
sonra <- c(12, 13, 18, 19, 24, 25)
wilcox.test(sonra, once, paired = TRUE,
            alternative = "two.sided", exact = TRUE,
            correct = FALSE)

# Uc bagimsiz grup
puan <- c(1,2,3, 4,5,6, 7,8,9)
grup <- factor(rep(c("A", "B", "C"), each = 3))
kruskal.test(puan ~ grup)
pairwise.wilcox.test(puan, grup,
                     p.adjust.method = "holm",
                     exact = TRUE)
\end{lstlisting}

\begin{outputguidebox}
\textbf{Çözüm ve kontrol değerleri.} Bağımsız iki grupta sıra toplamına
dayalı istatistik $W=25$ ve kesin iki yönlü $p=0{,}00794$'tür.
Eşleştirilmiş örnekte bütün farklar pozitif olduğundan $V=21$ ve kesin
$p=0{,}03125$ bulunur. Üç grup için $H(2)=7{,}20$, $p=0{,}02732$'dir.
Genel test hangi çiftlerin farklı olduğunu tek başına göstermez.
\end{outputguidebox}

\textbf{Yorum.} R, bağlı sıra veya sıfır fark bulunduğunda kesin yöntemi
kullanamayabilir ve uyarı verebilir. Bu durumda asimptotik yaklaşımın veya
Monte Carlo yönteminin kullanıldığı raporda belirtilmelidir. Mann--Whitney
sonucu, dağılım biçimleri için ek koşul olmadan otomatik ``medyan testi''
olarak sunulmamalıdır.

\textbf{Pekiştirme ve yanıt.} İkili karşılaştırmalarda
\texttt{p.adjust.method="holm"} aile içindeki çoklu test sorununu kontrol
eder. Düzeltilmemiş $p$-değerlerinden yalnız anlamlı olanları seçmeyin.